The concentration sensitivity of equiatomic TiNi, TiCu, and PdAu was examined at 298.15~K by varying one concentration at a time. The nominal aqueous metal-ion concentration was varied from $10^{-6}$ to $10^{-2}\,\mathrm{M}$, the total ammonia concentration from $0.002$ to $0.1\,\mathrm{M}$, and the total glycine concentration from $0.01$ to $1.0\,\mathrm{M}$. Other concentrations were held at their reference values: nominal aqueous metal-ion concentration $10^{-4}\,\mathrm{M}$, $[\ce{NH3}]_{\mathrm{tot}}=0.02\,\mathrm{M}$, and $[\mathrm{Gly}]_{\mathrm{tot}}=0.1\,\mathrm{M}$, without added cyanide. The nominal metal-ion concentration corresponds to the activity input under the ideal-solution convention. Each black outline marks pH 10.5--13.5 and $-2.0$ to $2.3$~V vs RHE, converted to SHE and clipped to the plotted range.

Only concentration variations that change the predicted phase map are shown below. Within the included thermodynamic model and the sampled grid resolution, the TiNi phase map is unchanged as the total ammonia concentration varies from $0.002$ to $0.1\,\mathrm{M}$, and the PdAu phase map is unchanged as the total glycine concentration varies from $0.01$ to $1.0\,\mathrm{M}$; these two triptychs are therefore omitted.

\begin{figure}[!htbp]
    \centering
    \includegraphics[width=\textwidth]{TiNi_activity_three_levels_1x3_side_colors.pdf}
    \caption{TiNi Pourbaix diagrams with the nominal aqueous metal-ion concentration varied from $10^{-6}$ to $10^{-2}\,\mathrm{M}$. Other concentrations were held at their reference values.}
    \label{fig:TiNi_activity_sensitivity_SI}
\end{figure}

\begin{figure}[!htbp]
    \centering
    \includegraphics[width=\textwidth]{TiNi_Gly_three_levels_1x3_side_colors.pdf}
    \caption{TiNi Pourbaix diagrams with the total glycine concentration varied from $0.01$ to $1.0\,\mathrm{M}$. Other concentrations were held at their reference values.}
    \label{fig:TiNi_Gly_sensitivity_SI}
\end{figure}

\begin{figure}[!htbp]
    \centering
    \includegraphics[width=\textwidth]{TiCu_activity_three_levels_1x3_side_colors.pdf}
    \caption{TiCu Pourbaix diagrams with the nominal aqueous metal-ion concentration varied from $10^{-6}$ to $10^{-2}\,\mathrm{M}$. Other concentrations were held at their reference values.}
    \label{fig:TiCu_activity_sensitivity_SI}
\end{figure}

\begin{figure}[!htbp]
    \centering
    \includegraphics[width=\textwidth]{TiCu_NH3_three_levels_1x3_side_colors.pdf}
    \caption{TiCu Pourbaix diagrams with the total ammonia concentration varied from $0.002$ to $0.1\,\mathrm{M}$. Other concentrations were held at their reference values.}
    \label{fig:TiCu_NH3_sensitivity_SI}
\end{figure}

\begin{figure}[!htbp]
    \centering
    \includegraphics[width=\textwidth]{TiCu_Gly_three_levels_1x3_side_colors.pdf}
    \caption{TiCu Pourbaix diagrams with the total glycine concentration varied from $0.01$ to $1.0\,\mathrm{M}$. Other concentrations were held at their reference values.}
    \label{fig:TiCu_Gly_sensitivity_SI}
\end{figure}

\begin{figure}[!htbp]
    \centering
    \includegraphics[width=\textwidth]{PdAu_activity_three_levels_1x3_side_colors.pdf}
    \caption{PdAu Pourbaix diagrams with the nominal aqueous metal-ion concentration varied from $10^{-6}$ to $10^{-2}\,\mathrm{M}$. Other concentrations were held at their reference values.}
    \label{fig:PdAu_activity_sensitivity_SI}
\end{figure}

\begin{figure}[!htbp]
    \centering
    \includegraphics[width=\textwidth]{PdAu_NH3_three_levels_1x3_side_colors.pdf}
    \caption{PdAu Pourbaix diagrams with the total ammonia concentration varied from $0.002$ to $0.1\,\mathrm{M}$. Other concentrations were held at their reference values.}
    \label{fig:PdAu_NH3_sensitivity_SI}
\end{figure}
